\newpage
\chapter{Course Registration and Participation Records}

\section{Preparing the Form and Response Sheet}

For the time being, course registration is accepted through a separate Google Form for each course that requires registration. The Dean responsible copies the Academy's template and, in consultation with the Lecturer, sets the course name, a display of or link to the syllabus, any conditions for taking the course, and contact details for changes and cancellation. Limit response fields to information necessary to manage participation.
Set a Google Sheets spreadsheet dedicated to that course as the destination for responses and give the Lecturer responsible the editing permissions needed to enter records and assessments. Before announcing that registration is open, check the form's display, the destination of responses, and the spreadsheet's sharing settings with the Lecturer.

\section{When Registration Takes Effect and Checking Conditions}

Course registration takes effect when a Student submits the required information through the course's form and the form accepts the response. If the Lecturer sets conditions for taking the course, display them in the form's instructions so that Students can check them before submitting, in accordance with Chapter 7, Article 7, paragraph 2 of the Academy Regulations.
The Lecturer checks registrations in the course's response sheet. If omissions, duplicate entries, or a failure to meet the conditions are found, check with the person concerned, explain how the matter will be handled, and record the action taken. When a Lecturer seeks advice, the Dean checks the registration record and the conditions that were communicated, and provides support. Do not treat a failure to meet course conditions as a loss of Student or Lecturer status.

\section{Changes, Cancellation, and Corrections}

Requests to change or cancel a registration go to the Lecturer responsible; tell Students how to contact that Lecturer. The Lecturer confirms the registration concerned and the person's request, records the status, such as ``taking the course'' or ``cancelled'', and the changes and other relevant details in the response sheet, and informs the person of the result.
Correct errors such as a mistaken Student ID after checking with the person concerned. If a Dean receives a request, pass it to the Lecturer and provide support as needed.

\section{Assessment Records and Permitted Uses}

The Lecturer may record participation and assessments within the course in its response sheet. Use records only to identify participants, change or cancel registrations, communicate as necessary to run the course, and enable the Lecturer responsible to record participation and assessments. Do not reuse them as lists for other courses or for purposes outside the Academy.
Assessment records kept by a Lecturer do not constitute the Academy's formal recognition of grades or credits. Full operation of the grades and credits system is planned for the next academic term or later.

\section{Retention, Permissions, and Handover}

Access to the response sheet is limited to the Lecturer responsible and Deans who need it for administrative work. The Dean responsible checks sharing settings and, when the course ends, reviews the need to retain records with the Lecturer and removes unnecessary copies and sharing permissions.
No uniform retention period or deletion schedule has yet been established. Before deleting individual records, discuss their handling among the Deans and record the decision.
When the Lecturer or Dean responsible changes, hand over pending changes and corrections, the location of records, and the necessary viewing and editing permissions to the successor. Remove permissions that are no longer needed after the handover.
